\subsection{Actual harmonic expressions}\label{subsec:fc-line-harmonic}

\begin{remark}[Touching holes and strong maxima]\label{rmk:fc-harmonic-touching}
\uses{def:fc-standard-affinoid}
\source{gr25-classically-psh:prop:harmonic.properties}
The strong maximum principle is local, and extends to a connected component
of the interior. It need not imply constancy on the whole closed affinoid
when distinct open holes touch at a Gauss boundary point. Indeed, choose
$0<r<|t|<s<R$ and set
$U=D(0,R)\setminus(D^-(0,|t|)\cup D^-(t,r))$.
The harmonic expression $\log|T-t|-\log|T|$ is at most zero on $U$ and
equals zero at the nonboundary skeleton point $\zeta_{0,s}$, but is strictly
negative at $\zeta_{t,r}$. The internal vertex $\zeta_{0,|t|}$ is a boundary
point. Thus the whole-affinoid conclusion requires a different hypothesis;
the local and interior-component formulations apply to the permitted data.
\end{remark}

\begin{definition}[Actual logarithmic harmonic functions]\label{def:fc-actual-harmonic}
\lean{EnvelopeProperty.Berkovich.RelativeSpectrum.FCStandardAffinoid.harmonicExpression, EnvelopeProperty.Berkovich.RelativeSpectrum.FCStandardAffinoid.IsHarmonic}
\leanok
\uses{def:fc-standard-affinoid}
For $c\in\R$ and $\gamma\in\R^m$, put
$H_{c,\gamma}(x)=c+\sum_{j=1}^m\gamma_j\log|T-a_j|_x$.
A real function on the actual subspace $U$ is harmonic if it equals the
restriction of some $H_{c,\gamma}$. Coefficients are witnesses for the
actual function; distinct witnesses may define the same function when
boundary points coincide.
\end{definition}

\begin{lemma}[Continuity and skeleton invariance]\label{lem:fc-harmonic-retraction}
\lean{EnvelopeProperty.Berkovich.RelativeSpectrum.FCStandardAffinoid.continuous_harmonicExpression, EnvelopeProperty.Berkovich.RelativeSpectrum.FCStandardAffinoid.IsHarmonic.continuous, EnvelopeProperty.Berkovich.RelativeSpectrum.FCStandardAffinoid.harmonicExpression_skeletonProjection, EnvelopeProperty.Berkovich.RelativeSpectrum.FCStandardAffinoid.IsHarmonic.retraction_invariant, EnvelopeProperty.Berkovich.RelativeSpectrum.FCStandardAffinoid.harmonicExpression_gauss}
\leanok
\uses{def:fc-actual-harmonic, thm:fc-radial-retraction}
Every such function is continuous on $U$ and satisfies
$H_{c,\gamma}(\tau_U(x))=H_{c,\gamma}(x)$.
On a radius segment its exact profile is
\[
 H_{c,\gamma}(\zeta_{b,s})=
 c+\sum_j\gamma_j\log\max\{|b-a_j|,s\}.
\]
\end{lemma}

\begin{proof}
On $U$, the $j$th linear value is at least the strictly positive hole radius
$r_j$. Its logarithm is therefore continuous; finite real linear combinations
remain continuous. The projection inflates by the capped minimum of the
hole-coordinate values. This radius is at most every hole-coordinate value,
so inflation leaves every such coordinate unchanged, and hence leaves the
displayed function unchanged. The segment formula follows by evaluating each
linear coordinate at the original Gauss point.
\end{proof}

\begin{lemma}[A constant piece away from finite Gauss centres]\label{lem:fc-harmonic-nongauss-constant}
\lean{EnvelopeProperty.Berkovich.RelativeSpectrum.FCStandardAffinoid.harmonicExpression_eventually_eq_of_min_nongauss}
\leanok
\uses{def:fc-actual-harmonic, lem:fc-strict-residue-disc}
At a point where all the finitely many hole coordinates are positive, choose
one having the smallest value $s$. If the point is not the Gauss point with
that hole centre and radius $s$, every harmonic expression is constant in
an actual open neighbourhood of the point.
\end{lemma}

\begin{proof}
Choose a strict residue disc $D^-(b,s)$ containing the point. Each original
hole coordinate has value at least $s$, whereas $|T-b|<s$. Unequal-value
ultrametric equality implies that each hole centre has distance from $b$
equal to the corresponding coordinate value at the point. These distances
are all at least $s$. At every point of $D^-(b,s)$, unequal-value equality
therefore makes each hole coordinate equal to that same constant distance.
All logarithms, and consequently their real linear combination, are constant
throughout this open disc.
\end{proof}

\begin{lemma}[Local constancy at actual Gauss maxima]\label{lem:fc-harmonic-gauss-maximum}
\lean{EnvelopeProperty.Berkovich.RelativeSpectrum.FCStandardAffinoid.harmonicExpression_eventually_eq_of_isLocalMax_gauss}
\leanok
\uses{def:fc-actual-harmonic, analytic:actual-original-unit-max-gauss-local-constancy}
At a positive-radius actual Gauss point, a harmonic expression having a
local maximum is constant in an actual open neighbourhood.
\end{lemma}

\begin{proof}
Translate the Gauss centre to zero. Each logarithm becomes the real logarithm
of the original rational unit $(T-(a_j-b))/1$ near the translated Gauss point;
these units are nonzero there. Use the actual rational-unit finite-family
collars and slopes. A local maximum makes every outward slope nonpositive.
The slopes balance, so all are zero. The actual finite listed-branch collars
and every unlisted residue branch then have the same function value. Their
constructed open neighbourhood contains the Gauss point. Translating back
gives the asserted open neighbourhood of the original point.
\end{proof}

\begin{theorem}[Local constancy at every actual maximum]\label{thm:fc-harmonic-local-maximum}
\lean{EnvelopeProperty.Berkovich.RelativeSpectrum.FCStandardAffinoid.harmonicExpression_eventually_eq_of_isLocalMax}
\leanok
\uses{lem:fc-harmonic-nongauss-constant, lem:fc-harmonic-gauss-maximum}
At a point where every hole-coordinate value is positive, an actual harmonic
expression having a local maximum is constant in an actual neighbourhood.
\end{theorem}

\begin{proof}
If there are no holes, the expression is constant. Otherwise choose a
smallest one of the finitely many positive coordinate values, say at site
$a_j$ with value $s$. If the point is $\zeta_{a_j,s}$, apply the positive-radius
Gauss maximum lemma. If it is different, the strict-residue constant-piece
lemma applies. These two cases cover every original point, including
radius-zero and nested-disc points.
\end{proof}
